\documentclass[preview]{standalone}

\usepackage{amsmath}
\usepackage{amssymb}
\usepackage{stellar}
\usepackage{definitions}

\begin{document}

\id{psicologia-condizionamento-operante}
\genpage

\section{Condizionamento operante}

\begin{snippetexample}{film-horror-condizionamento-operante-example}{Film horror e comportamento futuro}
    Dopo aver visto un film horror, una persona potrebbe decidere di non andare
    più al cinema a vedere film horror. In questo caso l'apprendimento non
    riguarda solo l'associazione tra due stimoli, ma il modo in cui una
    conseguenza modifica la probabilità di ripetere un comportamento.
\end{snippetexample}

\begin{snippetexample}{thorndike-gabbia-problema-example}{Gabbia-problema di Thorndike}
    Thorndike studiò gatti chiusi in una gabbia-problema. Inizialmente gli
    animali producevano molte azioni impulsive; quando una di esse permetteva di
    aprire la porta, le azioni inefficaci venivano progressivamente abbandonate
    e quella efficace diventava più frequente. L'apprendimento avveniva per
    prove ed errori, attraverso l'associazione tra la situazione e la risposta
    capace di produrre un esito soddisfacente.
\end{snippetexample}

\includesnpt[width=45\%,src=/snippet/static-psychology/thorndike-problem-box.jpg]{centered-img}

\begin{snippetdefinition}{legge-effetto-definition}{Legge dell'effetto}
    La \emph{legge dell'effetto} afferma che le risposte seguite da conseguenze
    soddisfacenti tendono a essere ripetute, mentre quelle prive di effetto o
    seguite da conseguenze sfavorevoli tendono a diminuire.
\end{snippetdefinition}

\begin{snippetdefinition}{condizionamento-operante-definition}{Condizionamento operante}
    Il \emph{condizionamento operante} è una procedura attraverso cui si
    manipolano le conseguenze del comportamento di un organismo per valutarne
    l'effetto sul comportamento successivo.
\end{snippetdefinition}

\begin{snippetdefinition}{comportamento-operante-definition}{Comportamento operante}
    Un \emph{comportamento operante} è qualunque comportamento messo in atto
    dall'organismo capace di produrre effetti osservabili sull'ambiente. A
    differenza dei riflessi del condizionamento classico, non è attivato
    automaticamente da uno stimolo incondizionato.
\end{snippetdefinition}

\begin{snippetexample}{skinner-box-example}{Skinner box}
    Skinner inventò la \emph{Skinner box}, o camera operante, per studiare come
    le conseguenze ambientali modificano la probabilità di un comportamento. Se
    un ratto preme una leva e riceve cibo, lo sperimentatore può misurare come
    il rinforzo modifica la frequenza della pressione della leva e quali
    condizioni rendono più probabile il comportamento desiderato.
\end{snippetexample}

\includesnpt[width=47\%,src=/snippet/static-psychology/skinner-box-operant-chamber.jpg]{centered-img}

\begin{snippetexample}{piccione-cibo-example}{Piccione e cibo}
    In un esperimento operante, la beccata di un piccione su un bottone può
    essere seguita dall'emissione di cibo. Se il cibo è contingente alla beccata
    e non ad altri comportamenti, la frequenza delle beccate tende ad aumentare.
\end{snippetexample}

\begin{snippetdefinition}{rinforzo-definition}{Rinforzo}
    Il \emph{rinforzo} è qualunque stimolo che, somministrato in modo
    contingente a una risposta, aumenta la probabilità che quella risposta si
    verifichi.
\end{snippetdefinition}

\begin{snippetdefinition}{rinforzo-positivo-definition}{Rinforzo positivo}
    Il \emph{rinforzo positivo} si verifica quando un comportamento è seguito
    dall'erogazione di uno stimolo piacevole, aumentando la probabilità che quel
    comportamento si ripresenti.
\end{snippetdefinition}

\begin{snippetdefinition}{rinforzo-negativo-definition}{Rinforzo negativo}
    Il \emph{rinforzo negativo} si verifica quando un comportamento è seguito
    dalla rimozione, riduzione o prevenzione di uno stimolo avversivo,
    aumentando la probabilità che quel comportamento si ripresenti.
\end{snippetdefinition}

\begin{snippet}{fuga-evitamento}
    Nel rinforzo negativo si distinguono due casi. Nel \emph{condizionamento di
    fuga}, l'organismo impara una risposta che permette di scappare da uno
    stimolo avversivo già presente. Nel \emph{condizionamento di evitamento},
    impara una risposta che permette di evitare lo stimolo avversivo prima che
    compaia.
\end{snippet}

\begin{snippetnote}{rinforzo-positivo-negativo-note}{Rinforzo positivo e rinforzo negativo}
    Rinforzo positivo e rinforzo negativo aumentano entrambi la probabilità
    della risposta che li precede. Il rinforzo positivo aggiunge uno stimolo
    piacevole; il rinforzo negativo rimuove, riduce o previene uno stimolo
    spiacevole.
\end{snippetnote}

\begin{snippetdefinition}{stimolo-punitivo-definition}{Stimolo punitivo}
    Uno \emph{stimolo punitivo} è uno stimolo che, somministrato in modo
    contingente a una risposta, diminuisce la probabilità che quella risposta si
    verifichi.
\end{snippetdefinition}

\begin{snippetdefinition}{punizione-positiva-definition}{Punizione positiva}
    La \emph{punizione positiva} si verifica quando un comportamento è seguito
    dalla somministrazione di uno stimolo spiacevole, riducendo la probabilità
    che quel comportamento si ripresenti.
\end{snippetdefinition}

\begin{snippetexample}{punizione-positiva-example}{Punizione positiva}
    Bere un caffè bollente produce dolore sulla punta della lingua. In seguito,
    la persona può imparare ad aspettare prima di bere caffè molto caldo.
\end{snippetexample}

\begin{snippetdefinition}{punizione-negativa-definition}{Punizione negativa}
    La \emph{punizione negativa} si verifica quando un comportamento è seguito
    dalla rimozione di uno stimolo piacevole, riducendo la probabilità che quel
    comportamento si ripresenti.
\end{snippetdefinition}

\begin{snippetexample}{punizione-negativa-example}{Punizione negativa}
    Se un genitore toglie la paghetta al figlio dopo che ha litigato con il
    fratellino, il bambino può imparare a non ripetere quel comportamento.
\end{snippetexample}

\begin{snippetdefinition}{stimoli-discriminativi-definition}{Stimoli discriminativi}
    Gli \emph{stimoli discriminativi} sono stimoli che precedono una particolare
    risposta e indicano il contesto in cui quel comportamento produrrà
    probabilmente una certa conseguenza.
\end{snippetdefinition}

\begin{snippetdefinition}{contingenza-tre-termini-definition}{Contingenza a tre termini}
    La \emph{contingenza a tre termini}, secondo Skinner, è la sequenza composta
    da \emph{stimolo discriminativo}, \emph{comportamento} e
    \emph{conseguenza}. Rappresenta un concetto centrale del comportamentismo.
\end{snippetdefinition}

\includesnpt[width=65\%,src=/snippet/static-psychology/three-term-contingency.jpg]{centered-img}

\begin{snippetdefinition}{rinforzo-primario-definition}{Rinforzo primario}
    Un \emph{rinforzo primario} è un rinforzo, come cibo o acqua, le cui
    proprietà rinforzanti sono biologicamente determinate.
\end{snippetdefinition}

\begin{snippetdefinition}{rinforzo-condizionato-definition}{Rinforzo condizionato}
    Un \emph{rinforzo condizionato} è uno stimolo inizialmente neutro che,
    associandosi a un rinforzo primario, diventa capace di funzionare come
    rinforzo per risposte operanti.
\end{snippetdefinition}

\begin{snippetexample}{rinforzo-parziale-skinner-example}{Skinner e il rinforzo parziale}
    Skinner osservò che, quando i ratti ricevevano cibo solo dopo un certo
    intervallo di tempo, mantenevano un tasso di risposta simile a quello
    ottenuto con il rinforzo continuo e, durante l'estinzione, continuavano a
    rispondere più a lungo. Questo portò allo studio sistematico degli schemi di
    rinforzo parziale.
\end{snippetexample}

\begin{snippetdefinition}{schema-rapporto-fisso-definition}{Schema a rapporto fisso}
    Nello \emph{schema a rapporto fisso}, o \emph{RF}, il rinforzo viene
    somministrato dopo che l'organismo ha prodotto un numero prestabilito di
    risposte.
\end{snippetdefinition}

\begin{snippetexample}{rapporto-fisso-example}{Rapporto fisso}
    Un venditore deve vendere un certo numero di prodotti prima di ricevere il
    compenso. Il rinforzo dipende quindi dal raggiungimento di un numero fisso
    di risposte.
\end{snippetexample}

\begin{snippetdefinition}{schema-rapporto-variabile-definition}{Schema a rapporto variabile}
    Nello \emph{schema a rapporto variabile}, o \emph{RV}, il numero medio di
    risposte tra un rinforzo e il successivo è predeterminato, ma il rinforzo
    può arrivare dopo un numero variabile di risposte.
\end{snippetdefinition}

\begin{snippetexample}{rapporto-variabile-example}{Rapporto variabile}
    Il gioco d'azzardo è governato da schemi a rapporto variabile. La risposta
    di inserire gettoni in una slot machine viene mantenuta dalla vincita, che
    arriva dopo un numero sconosciuto e variabile di risposte.
\end{snippetexample}

\includesnpt[width=74\%,src=/snippet/static-psychology/reinforcement-schedules-ratio.jpg]{centered-img}

\begin{snippetdefinition}{schema-intervallo-fisso-definition}{Schema a intervallo fisso}
    Nello \emph{schema a intervallo fisso}, o \emph{IF}, il rinforzo viene
    erogato in seguito alla prima risposta fornita dopo un intervallo di tempo
    costante e prestabilito.
\end{snippetdefinition}

\begin{snippetdefinition}{schema-intervallo-variabile-definition}{Schema a intervallo variabile}
    Nello \emph{schema a intervallo variabile}, o \emph{IV}, la media
    dell'intervallo tra i rinforzi è determinata, ma il rinforzo viene
    somministrato dopo intervalli variabili.
\end{snippetdefinition}

\includesnpt[width=58\%,src=/snippet/static-psychology/reinforcement-schedules-interval.jpg]{centered-img}

\begin{snippetdefinition}{modellamento-definition}{Modellamento}
    Il \emph{modellamento}, o \emph{shaping}, è un metodo di apprendimento in
    cui si rinforza ogni risposta che si avvicina progressivamente al
    comportamento desiderato, fino a ottenere la risposta finale.
\end{snippetdefinition}

\begin{snippetdefinition}{rinforzo-differenziale-definition}{Rinforzo differenziale}
    Il \emph{rinforzo differenziale} consiste nel rinforzare selettivamente le
    risposte che si avvicinano al comportamento-obiettivo, trascurando quelle
    meno adeguate.
\end{snippetdefinition}

\begin{snippetdefinition}{apprendimento-associativo-definition}{Apprendimento associativo}
    L'\emph{apprendimento associativo} è una teoria meccanicistica secondo cui
    l'apprendimento viene misurato come variazione delle risposte
    comportamentali dopo una situazione di stimolazione.
\end{snippetdefinition}

\end{document}
